\subsection{مبانی نظری و فرمول‌بندی ریاضی}
فرض کنید ماتریس وزن لایه خطی با $\mathbf{W} \in \mathbb{R}^{d_{\text{out}} \times d_{\text{in}}}$ مشخص شود. در روش تفکیک علامت و بزرگی، ماتریس به دو بخش ماتریس علامت $\mathbf{S}_{\mathbf{W}} = \operatorname{sgn}(\mathbf{W})$ و ماتریس قدرمطلق غیرمنفی $\mathbf{A} = |\mathbf{W}| \ge 0$ تفکیک می‌گردد. هدف \lr{NMF} تقریب ماتریس $\mathbf{A}$ با دو ماتریس کم‌مرتبه غیرمنفی $\mathbf{H} \in \mathbb{R}_{+}^{d_{\text{out}} \times r}$ و $\mathbf{V} \in \mathbb{R}_{+}^{r \times d_{\text{in}}}$ است:
\begin{equation}
\min_{\mathbf{H} \ge 0, \mathbf{V} \ge 0} \|\mathbf{A} - \mathbf{H}\mathbf{V}\|_F^2
\end{equation}
برای دستیابی به همگرایی سریع و پایدار، از الگوریتم اولیه‌سازی \lr{NNDSVD (Non-negative Double Singular Value Decomposition)} استفاده می‌شود. با محاسبه تجزیه کم‌مرتبه $\mathbf{A} \approx \sum_{j=1}^r \sigma_j \mathbf{u}_j \mathbf{v}_j^\top$، مؤلفه‌های مثبت و منفی بردارهای منفرد تفکیک می‌گردند:
\begin{equation}
\mathbf{u}_j^+ = \max(\mathbf{u}_j, 0), \quad \mathbf{u}_j^- = \max(-\mathbf{u}_j, 0)
\end{equation}
ستون‌های اولیه $\mathbf{H}$ و سطرهای $\mathbf{V}$ بر اساس بزرگترین حاصل‌ضرب نرم برگزیده می‌شوند. در ادامه، قواعد به‌روزرسانی ضربی (\lr{Lee \& Seung}) روی کارت گرافیک اجرا می‌گردند:
\begin{equation}
\mathbf{V} \leftarrow \mathbf{V} \odot \frac{\mathbf{H}^\top \mathbf{A}}{\mathbf{H}^\top \mathbf{H} \mathbf{V} + \epsilon}, \quad \mathbf{H} \leftarrow \mathbf{H} \odot \frac{\mathbf{A} \mathbf{V}^\top}{\mathbf{H} \mathbf{V} \mathbf{V}^\top + \epsilon}
\end{equation}
در نهایت، ماتریس وزن بازسازی‌شده با اعمال مجدد ماتریس علامت حاصل می‌شود:
\begin{equation}
\hat{\mathbf{W}} = (\mathbf{H}\mathbf{V}) \odot \mathbf{S}_{\mathbf{W}}
\end{equation}


\paragraph{تحلیل طیفی و نامساوی‌های پایداری خطا}
برای بررسی میزان دقت تقریب ماتریس وزن یا نگاشت توجه در ماژول \lr{nmf\_compressor}، تفاضل ماتریس فاکتورگیری‌شده از نگاشت متراکم اصلی بر اساس نامساوی هافدینگ-فون‌نیومان ارزیابی می‌گردد. فرض کنید $\mathbf{W} \in \mathbb{R}^{d_{\text{in}} \times d_{\text{out}}}$ ماتریس وزن هدف باشد که توسط تقریب کم‌مرتبه یا پراکنده با رتبه موثر $r$ جایگزین شده است. کران بالای خطای بازسازی بر مبنای نرم فروبنیوس به صورت زیر تصریح می‌شود:
\begin{equation}
\label{eq:spectral_bound_nmf_compressor}
\| \mathbf{W} - \mathbf{W}_{r} \|_F = \sqrt{\sum_{i=r+1}^{\min(d_{\text{in}}, d_{\text{out}})} \sigma_i^2(\mathbf{W})} \le \epsilon \| \mathbf{W} \|_F
\end{equation}
که در آن $\sigma_i(\mathbf{W})$ بیانگر $i$-امین مقدار منفرد مرتب‌شده ماتریس به صورت نزولی است. نسبت حفظ انرژی طیفی $\eta_r = \frac{\sum_{i=1}^r \sigma_i^2}{\sum_{i=1}^{d} \sigma_i^2}$ تضمین می‌کند که با انتخاب آستانه $\tau \ge 0.90$، بیش از نود درصد از پراکندگی واریانس داده‌ها حفظ شده و افت سیگنال در انتشار گرادیان رو به عقب به حداقل ممکن کاهش می‌یابد.

\paragraph{تحلیل مجانبی و پیچیدگی اشتراک حافظه سخت‌افزاری}
از منظر پیچیدگی محاسباتی و انتقال داده‌ها در سلسومراتب حافظه (حافظه سراسری \lr{HBM} به حافظه محلی \lr{SRAM})، عملیات ضرب ماتریسی در این مکانیزم نیازمند تخصیص بافر موقت با ابعاد فاکتورگیری‌شده است. برای پردازش دنباله‌ای به طول $N$ با بعد مخفی $D$، بار محاسباتی هسته از فرم درجه دو به رابطه زیر تقلیل می‌یابد:
\begin{equation}
\label{eq:complexity_nmf_compressor}
\mathcal{C}_{\text{nmf\_compressor}} = \mathcal{O}\left( B \cdot N \cdot D \cdot r + B \cdot H \cdot N \cdot d_k \right) \ll \mathcal{O}\left( B \cdot H \cdot N^2 \cdot d_k \right)
\end{equation}
این بهبود مرتبه پیچیدگی، با کاهش چشمگیر حجم تبادل تانسورها در گذرگاه‌های ارتباطی کارت گرافیک، مانع از اشباع پهنای باند حافظه (پدیده \lr{Memory-Bound}) در فرآیند استنتاج طولانی‌مدت می‌گردد.
